\documentclass[10pt,a4paper]{amsart}
\usepackage[margin=19mm]{geometry}
\usepackage{amsmath,amssymb,mathtools}
\usepackage[scr=rsfs,cal=euler]{mathalfa}
\usepackage{xfrac}
\usepackage[unicode,hidelinks,bookmarks=false]{hyperref}
\newcommand{\ie}{i.e.\ }
\newcommand{\cf}{cf.\ }
\newcommand{\resp}{respectively\ }
\newcommand{\Subsection}{\subsection}
\providecommand{\nobreakdash}{\nobreak\hbox{-}}
\setlength{\emergencystretch}{2em}
\multlinegap0pt
\usepackage{bm}
\usepackage{bbm}
\usepackage{dsfont}
\usepackage{xspace}
\usepackage{cancel}
\usepackage{mathtools}
\usepackage[shortlabels]{enumitem}
\usepackage{amsthm}
\makeatletter
\@ifundefined{theorem}{\newtheorem{theorem}{Theorem}}{}
\@ifundefined{assumption}{\newtheorem{assumption}[theorem]{Assumption}}{}
\@ifundefined{definition}{\newtheorem{definition}[theorem]{Definition}}{}
\@ifundefined{lemma}{\newtheorem{lemma}[theorem]{Lemma}}{}
\makeatother
\newcommand{\C}{\mathcal C}
\newcommand{\E}{\mathbb{E}}
\newcommand{\F}{F}
\newcommand{\N}{\mathbb{N}}
\newcommand{\R}{\mathbb{R}}
\newcommand{\av}[1]{\left<#1\right>}
\newcommand{\cM}{\mathcal M}
\newcommand{\dd}{\text{d}}
\title[Well-Posedness for Dean-Kawasaki Models of Vlasov-Fokker-Pla]{Well-Posedness for Dean-Kawasaki Models of Vlasov-Fokker-Planck Type}
\author{Fenna M\"uller, Max von Renesse, Johannes Zimmer}
\date{Source: arXiv:2411.14334v3 (CC BY 4.0)}
\begin{document}
\maketitle

\noindent\emph{Source and licence.} Well-Posedness for Dean-Kawasaki Models of Vlasov-Fokker-Planck Type. Version arXiv:2411.14334v3 (CC BY 4.0), distributed under the Creative Commons Attribution 4.0 International licence, as recorded in the arXiv licence metadata for this exact version and shown on the abstract page https://arxiv.org/abs/2411.14334v3. Full rights notice: licenses/arxiv-2411.14334-CC-BY-4.0.txt.\par\smallskip

\noindent\textit{Editorial scope.} Counted unit: the Lemma whose source label is \texttt{lemma:secondmoment} in the article (source0.tex lines 1040--1046, source label \texttt{lemma:secondmoment}), delivered together with the article's own complete original proof of it (source0.tex lines 1048--1086, one \texttt{proof} environment of 39 lines). No mathematics is written, repaired or re-derived: statement and proof are the article's own text line for line. Required context retained uncounted: eqn:DKE1 (lines 374--377); assL:chainr (lines 481--497); assumption:interaction (lines 541--559); interact.L (lines 543--547); def:c2b\_etc (lines 1174--1188). Every definition, display and label of those lines is retained unchanged. Omitted from this package and not redistributed: 1--373, 378--480, 498--540, 548--1039, 1047--1047, 1087--1173, 1189--1264. Nothing inside a retained excerpt is omitted or altered: every retained body is delivered exactly as the article's lines are written, so no omission receipt of any kind accompanies this package. Citations: the retained excerpts carry no citation command at all, so no bibliography entry of the article is transcribed and no reference is redistributed. Reference adaptation: none was needed and none was made; every reference inside the delivered unit resolves inside it, so no unresolved reference or citation is produced. Printed slips kept literal: no printed word, symbol, hypothesis or number of the counted statement or of its proof was altered. Layout and class: the article's own class is \texttt{article} with packages including babel, inputenc, biblatex, csquotes, todonotes, subcaption, ulem, amssymb, amsmath, amsthm, exscale, mathtools, enumitem; the wrapper is the lane's pinned \texttt{amsart} editorial wrapper at 10pt on A4, which restates the theorem-like environments the retained text needs and adds the article's own macro definitions that the retained text needs (\texttt{\textbackslash C}, \texttt{\textbackslash E}, \texttt{\textbackslash F}, \texttt{\textbackslash N}, \texttt{\textbackslash R}, \texttt{\textbackslash av}, \texttt{\textbackslash cM}, \texttt{\textbackslash dd}), with extra wrapper packages (bm, bbm, dsfont, xspace, cancel, mathtools, enumitem). The delivered numbering is the unit's own. Assets: no retained span contains a figure environment, a graphics-inclusion command, a PDF-page inclusion, a table or a TikZ picture; no image file is copied into this package, the package declares no asset dependency and no image file is redistributed with it. Licence: version arXiv:2411.14334v3 of this article is distributed under the Creative Commons Attribution 4.0 International licence, as recorded in the arXiv licence metadata for this exact version and shown on the abstract page https://arxiv.org/abs/2411.14334v3. A full rights notice is delivered at licenses/arxiv-2411.14334-CC-BY-4.0.txt. Characters: every retained excerpt is pure ASCII, so the delivered engine, XeLaTeX with its default Latin Modern text font, reports no missing character.
\begin{align}
    \partial_t \mu_t = \alpha L^* \mu_t + \alpha \nabla \cdot(\mu_t \F
    _{\mu_t}) +  \nabla \cdot(\sqrt{\mu_t} \sigma \dot{W}_{z,t}), \label{eqn:DKE1}
\end{align}

  \begin{enumerate}[label=(L.{\arabic*})]
  \item \label{assL:martp}
    The operator $L$ acting on $\C_b^2(\R^k)$ admits a unique in law Markovian family of diffusion processes
    $X=(X_t^z)_{t \geq 0}^{z\in \R^k}$ solving the associated
    $(L,\C_b^2(\R^k), \delta_z)_{z \in \R^k}$-martingale problems.
       
  \item \label{assL:chainr} There is a set $\mathcal D\subset \C^2_b(\R^k) $ which is dense with respect to
    the topology of locally uniform convergence and is stable under composition with functions
    $\psi \in \C^\infty(\R)$ satisfying $\psi(0) = 0$ such that for all $\varphi \in \mathcal D$ and
    $t \geq 0$, it holds that $P_t \varphi \in \C^2_b(\R^k)$ and $P_t L\varphi = LP_t \varphi \in \C_b(\R^k)$,
    where $P_t \varphi(z) = \E \varphi(X_t^z)$.
  \item \label{assL:gradbound} For all $T>0$ and every function $\varphi \in \mathcal D$, the function
    $(z,t) \to \bigl|\sigma^T\nabla (P_t \varphi)\bigr|(z)$ is uniformly bounded on $\R^k\times [0,T]$.
  \item \label{assL:exhaust} There exists an exhaustion, i.e., a sequence of monotonically increasing sets $(A_n)_{n \in \N} \subset \mathcal{B}(\R^k)$ with
    $A_n \nearrow \R^k$, such that for any $t >0$ there exists a sequence $(c_n)_{n \in \N} \subset [0, 1)$
    such that $P_t 1_{A_n}(z) \leq c_n$ for all $z \in \R^k$ and $n \in \N$.
  \end{enumerate}

\begin{assumption}
\label{assumption:interaction} 
 \begin{enumerate}[label=(L.{\arabic*})] \addtocounter{enumi}{4}
 \item \label{interact.L} The fields $\sigma$ respectively $b$ are continuous and bounded respectively of at
   most linear growth. That is, there is some $K>0$ such that $|\sigma(z)|\leq K$ and $|b(z)|\leq K(1+|z|)$
   for $z \in \R^k$.
\end{enumerate}
\begin{enumerate}[label=(F.{\arabic*})]
\item \label{F-F1} %
  For some $G \in \C^2_{loc}(\cM_1^2)$ as defined in Definition~\ref{def:c2b_etc}, we have
  \begin{equation*}
    F_\mu(\cdot) = \left(\sigma \sigma^T\nabla\frac {\delta G }{\delta \mu}\right)(\cdot).
  \end{equation*}
\item \label{F-F2} It holds that $\sigma^ T\nabla \frac {\delta G}{\delta \mu}(\mu,.)$ is bounded and
  $\nabla^l \frac{\delta^m G}{\delta \mu^m}(\mu,.)$ is for
  $m,l\in \{1,2\}$ of at most linear growth, locally uniformly with respect to $\mu$ on $\cM_1$ and
  $\cM^2_1$ respectively.
\end{enumerate}
\end{assumption}

\begin{definition}
  \label{def:c2b_etc}
  A function $E\colon \cM^2 \mapsto \R$ belongs to the class $\C^2_{loc}(\cM^2)$ if
  \begin{itemize}
  \item $\frac{\delta E}{\delta \mu}(\mu,.) \in \C^2(\R^k) $ and
    $\frac{\delta^ 2 E}{\delta \mu^2}(\mu,.,.) \in \C^2(\R^k\times\R^k) $ exist for all $\mu \in \cM^2_1$;
  \item $E(.)$, $\frac{\delta E}{\delta \mu}$, and $\frac{\delta^ 2 E}{\delta \mu^2}$ as well as their first
    and second order spatial derivatives $\nabla^{l}\frac{\delta E}{\delta \mu}(.,)$ for $l\in \{1,2\}$ are
    (jointly) continuous on $\cM^2_{1,c}$, $\cM^2_{1,c}\times \R^k$ respectively $\cM^2_{1,c}\times \R^k\times \R^k$.
  \end{itemize}
  We say that a function $H\colon \cM^2_1\times \R^n \to \R $ is of \emph{at most linear growth, locally
    uniformly} on $\cM^2_1$, if for all $c>0$ there is a constant $K>0$ such that
  \[ H(\mu, z) \leq K (1+ |z|) \mbox{ for all } (\mu,z) \in \cM^2_{1,c} \times \R^n.\]
  $H$  is called \emph{bounded uniformly} on $\cM_1$ if  $\sup_{\mu \in \mathcal M_1, z \in \R^k } |H(\mu,z)| < \infty$. 
\end{definition}

\begin{lemma}
  \label{lemma:secondmoment}
  Let $b$ be of at most linear growth and $\F\colon\mathcal M_1 \times \R^k \to \R^k$ be continuous with
  $\sup_{\mu \in \mathcal M_1, z \in \R^k} |H(\mu,z)| < \infty$.  Assume that $(\mu_t)_{t \geq 0}$ is a
  solution to Equation~\eqref{eqn:DKE1} with initial condition $\mu_0 \in \cM^2_1$. Then $\mu_t \in \cM^2_1$
  almost surely for all $t \geq 0$.
\end{lemma}

\begin{proof} The choice of the test function $\varphi =1$ in the weak martingale formulation of
  \eqref{eqn:DKE1} reveals that the total mass $\mu_t(\R^k)$ is conserved, with value $1$. In particular, by
  the assumptions on $F$, there is some constant $c$ such that $\sup_{t\geq 0, z \in \R^k} |F(\mu_t,z)| <c$
  almost surely. Here and in below we use the symbol $c$ for a generic constant whose value may change from
  line to line. Next, for $R > 0$, choose $\psi_R \in \C^\infty_b(\R)$ a smooth approximation of the function
  $ r \mapsto (r+1) \wedge R$ on $\R_{\geq 0}$ such that $\psi_R(r) = R$ for $r \geq R$, $r \leq \psi_R(r)$
  for all $r \leq R$, and $ \psi_R' \leq 1$ and $ |\psi_R''| \leq c $ for uniformly in $R$.  Then
  $\varphi_R = \psi_R \circ e_2 \in \C^2_b(\R^k)$ and the chain rule for $L$ (cf.\
  Condition~\ref{assL:chainr}) together with~\ref{interact.L} of Assumption~\ref{assumption:interaction}
  yields
  \begin{equation*}|L \varphi_R(z)| \leq c \left(1 +\varphi_R(z)\right),\end{equation*}
  with a constant $c$ independent of $R$. 
  Now, since $\varphi_R \in \C^2_b(\R^k)$, by definition of $\mu_t$ the process $M_t(\varphi_R)$ is a
  continuous martingale and hence $\tau_{N,R} := \inf\{t \geq 0 \bigm|\av{\mu_t, \varphi_R}\geq N\}$ is a
  well-defined stopping time, with $\lim_{N \to \infty}\tau_{N,R} = \infty$ almost surely, for all $R >
  0$. Therefore,
 \begin{align*}
                \E\left( \av{\mu_{t\wedge\tau_N}, \varphi_R} \right) &= \E\left(\av{\mu_0, \varphi_R} + \int_0^{t\wedge\tau_N} \av{\mu_s, \alpha L\varphi_R+F_{\mu_s} \cdot \nabla \varphi_R}\dd s\right)\\
                &\leq \E\left( \av{\mu_0, \varphi_R} \right) + \int_0^t \E\left( \av{\mu_{s\wedge \tau_{N,R}}, c\left(1+ \varphi_R\right)} \right)\dd s,
\end{align*}
Since up to the stopping time $\tau_{N,R}$, the process is bounded,
$\E\left( \av{\mu_{s\wedge \tau_{N,R}},\varphi_R} \right)$ is a continuous function. Hence, by Gr\"onwall's
Lemma
\begin{align*}
  \E\left( \av{\mu_{t \wedge \tau_{N,R}}, \varphi_R} \right) &\leq \left( \av{\mu_0, \varphi_R} + \av{\mu_0}ct\right)e^{ct},
\end{align*}
where we use again that the mass is preserved almost surely.  Using Fatou's lemma again, we finally
obtain
\begin{align*}
  \E\left( \av{\mu_{t},\varphi_R} \right) &= \E\left(\liminf_{R} \av{\mu_{t \wedge \tau_{N,R}},\varphi_R} \right)
                                            \leq \left( \av{\mu_0, \varphi_R} + ct\right)e^{ct}.
\end{align*}
{Due to $e_2(z) = |z|^2 \leq \lim_{R \to \infty}\varphi_R$,  Fatou's Lemma yields}
\begin{align}
  \label{est:secmombound}
  \E\left( \av{\mu_t, e_2} \right) &\leq \left( \av{\mu_0, 1+ e_2} + ct\right)e^{ct}
\end{align} 
and so in particular $\av{\mu_t, e_2} < \infty$ almost surely for all $t \geq 0$.
\end{proof}

\end{document}
